\section{The Problem Landscape}
\label{sec:landscape}

\subsection{The Core Phenomenon}

Many failures of initiative have the same business signature: the agent has a
declared long-term objective, a known next step, and sufficient conceptual
agreement with the objective, yet still does not begin. The relevant puzzle is
not ignorance. It is transition failure.

\begin{definition}[Behavioural State]
A \textbf{behavioural state} $s \in \StateSpace$ is a persistent bundle of
activity, posture, tools, affective expectation, and environmental cues.
\end{definition}

\begin{definition}[Target State]
A \textbf{target state} $s^\star$ is a behavioural state whose long-run value
dominates the current state under the agent's declared objective.
\end{definition}

\subsection{Three Observable Failure Modes}

\begin{enumerate}
  \item \textbf{Reward lock-in.} The present state pays immediately, even if it
        underperforms strategically.
  \item \textbf{Switching ambiguity.} The reader knows the target activity in
        general but does not know the exact first move.
  \item \textbf{Environmental capture.} The room, posture, devices, and nearby
        alternatives continuously reinforce the current state.
\end{enumerate}

\begin{examplebox}
\textbf{Study-start failure.} A student intends to read a chapter in algebra.
The desk is cluttered, the phone is open, the textbook is not prepared, and
there is no marked stopping point from the previous session. The student is not
lacking motivation in the abstract. The transition path is underspecified and
the current environment is paying more immediate reward than the target state.
\end{examplebox}

\subsection{Why ``Try Harder'' Is A Weak Instruction}

\begin{proposition}[Grit Alone Is Operationally Underspecified]
If a transition rule depends on repeated high-effort override with no reduction
in switching cost, then the rule is fragile under fatigue, interruption, and
environmental reinforcement.
\end{proposition}

\begin{discussion}
The proposition does not deny that effort matters. It states that effort is a
poor primary control variable. A business process that works only when the
operator happens to feel unusually strong is not a robust process. By the same
logic, a personal productivity system that depends on repeated heroic effort is
structurally weak.
\end{discussion}

\subsection{A Diagnostic Vocabulary}

For the rest of this manuscript, failures are classified by the dominant
constraint:
\begin{itemize}
  \item \textbf{Reward-dominant failure:} the current state pays too well.
  \item \textbf{Context-dominant failure:} the next action is too vague.
  \item \textbf{Environment-dominant failure:} the room or device ecosystem
        keeps reconstructing the old state.
  \item \textbf{Timing-dominant failure:} the agent attempts major transition
        when volitional capacity is already depleted.
\end{itemize}

\begin{summarybox}
\begin{itemize}
  \item The central problem is not moral weakness but transition design.
  \item Richer diagnosis requires naming \emph{which} layer blocks movement.
  \item This prepares the shift from qualitative complaint to formal model.
\end{itemize}
\end{summarybox}

